\documentclass[11pt]{article}

\usepackage[a4paper,margin=1in]{geometry}
\usepackage{amsmath,amssymb}
\usepackage{hyperref}

\hypersetup{
    colorlinks=true,
    linkcolor=blue,
    citecolor=blue,
    urlcolor=blue
}

\title{Gradient Descent on a Riemannian Manifold}
\author{}
\date{}

\begin{document}

\maketitle

\section*{The Euclidean Picture}

Let
\[
f:\mathbb R^n\to\mathbb R
\]
be smooth. In Euclidean space, the gradient $\nabla f(x)$ is characterized by
\[
Df(x)[v]
=
\nabla f(x)\cdot v
\]
for every direction $v\in\mathbb R^n$.

Here $Df(x)[v]$ is the
\href{directional_derivative.html}{directional derivative} of $f$ at $x$ in
the direction $v$. The dot product is the Euclidean metric. Thus the gradient
is not merely ``the list of partial derivatives.'' It is the vector that
\href{gradient_represents_derivative.html}{represents the linear map $Df(x)$
through the chosen inner product}.

\section*{Differential Versus Gradient}

On a smooth manifold $M$, the derivative of a smooth function
\[
f:M\to\mathbb R
\]
at a point $p\in M$ is naturally a covector
\[
df_p:T_pM\to\mathbb R.
\]
It eats a tangent vector $X_p\in T_pM$ and returns the directional derivative
\[
df_p(X_p)=X_p[f].
\]

So the differential $df_p$ lives in the cotangent space $T_p^*M$, not in the
tangent space $T_pM$. To turn this covector into a vector, we need
\href{metric_needed_for_gradient.html}{extra geometric structure}.

\section*{The Metric Converts Covectors Into Vectors}

A Riemannian metric assigns to each point $p\in M$ an inner product
\[
g_p:T_pM\times T_pM\to\mathbb R.
\]
The Riemannian gradient of $f$ is the unique tangent vector
\[
\operatorname{grad}_g f(p)\in T_pM
\]
such that
\[
df_p(X_p)
=
g_p(\operatorname{grad}_g f(p),X_p)
\]
for every tangent vector $X_p\in T_pM$.

This is the precise meaning of the phrase:
\begin{quote}
\href{gradient_represents_derivative.html}{The gradient is the vector that
represents the directional derivative through the metric}.
\end{quote}

The directional derivative is the covector $df_p$. The metric $g_p$ converts
that covector into the vector $\operatorname{grad}_g f(p)$.

\section*{Coordinate Formula}

Let $(x^1,\ldots,x^n)$ be local coordinates. Write the metric as
\[
g
=
g_{ij}\,dx^i\otimes dx^j.
\]
The differential is
\[
df
=
\frac{\partial f}{\partial x^i}\,dx^i.
\]
If
\[
\operatorname{grad}_g f
=
(\operatorname{grad}_g f)^i\frac{\partial}{\partial x^i},
\]
then the defining identity
\[
df(X)=g(\operatorname{grad}_g f,X)
\]
implies
\[
\frac{\partial f}{\partial x^j}
=
g_{ij}(\operatorname{grad}_g f)^i.
\]
Multiplying by the inverse metric $g^{kj}$ gives
\[
(\operatorname{grad}_g f)^k
=
g^{kj}\frac{\partial f}{\partial x^j}.
\]
Thus
\[
\operatorname{grad}_g f
=
g^{ij}\frac{\partial f}{\partial x^j}
\frac{\partial}{\partial x^i}.
\]

When $g_{ij}=\delta_{ij}$, this reduces to the usual Euclidean gradient.

\section*{Gradient Flow}

The Riemannian gradient flow of $f$ is the curve $p_t\in M$ solving
\[
\dot p_t
=
-\operatorname{grad}_g f(p_t).
\]
This is the intrinsic version of gradient descent. The negative sign means
that the curve moves in the
\href{steepest_descent_metric.html}{direction of steepest decrease as measured
by the metric $g$}.

Indeed, along a solution,
\[
\frac{d}{dt}f(p_t)
=
df_{p_t}(\dot p_t).
\]
Using the defining property of the gradient,
\[
df_{p_t}(\dot p_t)
=
g_{p_t}(\operatorname{grad}_g f(p_t),\dot p_t).
\]
Substituting
\[
\dot p_t=-\operatorname{grad}_g f(p_t)
\]
gives
\[
\frac{d}{dt}f(p_t)
=
-g_{p_t}(\operatorname{grad}_g f(p_t),\operatorname{grad}_g f(p_t))
\le 0.
\]
So the objective decreases monotonically.

\section*{Why This Matters for Wasserstein Geometry}

The Wasserstein calculation follows the same pattern.

On a finite-dimensional Riemannian manifold:
\[
df_p(X_p)
=
g_p(\operatorname{grad}_g f(p),X_p).
\]
In Wasserstein geometry, a tangent direction is represented by a velocity field
$v$, and the metric is the kinetic-energy inner product
\[
\langle u,v\rangle_\rho
=
\int_{\mathbb R^d}u(x)\cdot v(x)\,\rho(x)\,dx.
\]
Thus the Wasserstein gradient is identified by writing the directional
derivative in the form
\[
D\mathcal F(\rho)[v]
=
\langle \operatorname{grad}_{W_2}\mathcal F(\rho),v\rangle_\rho.
\]
This is why, once we obtain
\[
D\mathcal F(\rho)[v]
=
\int_{\mathbb R^d}
\nabla\left(\frac{\delta\mathcal F}{\delta\rho}\right)
\cdot v\,\rho\,dx,
\]
we identify the velocity-field representative of the Wasserstein gradient as
\[
\operatorname{grad}_{W_2}\mathcal F(\rho)
=
\nabla\left(\frac{\delta\mathcal F}{\delta\rho}\right).
\]

\end{document}
